\chapter{A Software Modem: DSC and NAVTEX}
\label{ch:softmodem}\index{FSK}\index{DSC}\index{NAVTEX}\index{SITOR-B}

Maritime radio still runs on two data systems designed in the 1970s and alive
in every GMDSS installation today. Digital Selective Calling (DSC, ITU-R
M.493) is the calling channel: a short FSK burst that addresses one station,
carries a category and telecommands, and ends in a checksum. NAVTEX (ITU-R
M.540, carried by SITOR-B, ITU-R M.476) is the broadcast channel: navigational
warnings as forward-error-corrected telex on 518\,kHz. Both are 100 baud FSK
with a 170\,Hz shift. Both repeat every character for time diversity. And both
are small enough to implement completely --- transmitter, receiver, and the
signal processing between --- in pure Ada on this SoC, with no DSP library, no
lookup tables, and no recorded audio.

This chapter is a case study rather than a driver manual: the stack lives in
the \texttt{esp32s3\_speak} example (packages \texttt{FSK}, \texttt{DSC} and
\texttt{Navtex}, each with a \texttt{Decoder} child), built on the full-duplex
I2S streaming of Chapter~\ref{ch:hal2}. What makes it worth a chapter is the
shape of the verification: every layer is checked on the host, the whole chain
is checked \emph{through the air} --- the board broadcasts a message from its
speaker and decodes what its own microphone hears --- and the protocol logic
is formally proved free of run-time errors.

\section{The modem}

One package, \texttt{FSK}, owns the physical layer for both protocols. The
transmitter is simple: collect bits, render each as one of two tones,
phase-continuous, with a click-free fade at the ends. The receiver is where
the design lives.

\textbf{Tone detection.} Each tone gets a narrow complex one-pole resonator:
rotate the running state by the tone frequency, shrink it by a radius, add the
sample. The decision is the difference of the two resonator energies. This
rejects noise and the opposite tone inherently. (A quadrature discriminator
was tried first: its mixing image passed the weak one-pole filter and flipped
the decision every sample.)

\textbf{Carrier offset.} A slowly tracked bias on the decision absorbs a
mistuned carrier. Slow is the design point: with a five-bit time constant the
bias chased any long run of one tone toward zero and wrecked the timing; at
thirty bits it removes only the static offset. Measured tolerance is
$\pm$80\,Hz --- nearly half the shift.

\textbf{Bit timing.} A transition-tracking DPLL: the bit phase advances every
sample, and a data edge pulls it toward mid-bit, where the decision is
sampled. It follows a $\pm$3\% baud error. Two later hardenings matter for
real signals. Corrections are \emph{gated} by the recent tone separation --- a
quarter-bit average, because the edge sample itself is a zero crossing and
always looks unconfident --- so a noise burst leaves the clock coasting on its
own period instead of random-walking it off frequency. And each edge may move
the phase at most twelve samples, so a burst's few spurious edges can never
add up to a bit slip.

\textbf{Soft decisions.} Alongside each hard bit the demodulator reports a
confidence: the tone separation at the sampling instant, the decision over the
total tone energy. The ratio is dimensionless, so it is independent of signal
level. A slow average of it is exported as a 0\,..\,255 link-quality meter:
255 on the board's own acoustic loopback, near zero on noise. (One instructive
dead end: an integrate-and-dump decision is textbook, but this DPLL dumps at
\emph{mid}-bit, so the integral straddles two bits and cancels. The
instantaneous mid-bit read is the aligned one.)

\section{The protocol layers}

DSC and NAVTEX share their skeleton, inherited from SITOR. Characters carry
their own check: DSC's 10-unit code is seven information bits plus a 3-bit
count of the zeros; NAVTEX's CCIR~476 code is seven bits of which exactly four
are marks. Each character is transmitted twice --- the DX copy, then four
other characters, then the RX repeat, 280\,ms of time diversity against burst
interference. A phasing preamble lets the receiver lock both the character
framing and the DX/RX slot parity.

The receive rule is the specification's, extended one step. Print an
unmutilated DX copy; failing that, the RX repeat; failing \emph{that} ---
where the specification prints an error mark --- rebuild the character bit by
bit from the two soft copies, letting whichever copy received each bit more
cleanly win, and accept the rebuild only if its check passes. On the test
harness, damage split across the two copies (each fails its check alone) now
decodes cleanly; only a character whose bits are lost in \emph{both} copies is
beyond recovery, and the checksum still catches it.

Polarity solves itself, by a pleasant algebra. The complement of a valid
10-unit character is also valid, with symbol $127-S$; the complement of a
CCIR~476 character has three marks instead of four and is invalid. So both
decoders hunt the phasing in true and inverted form, note which matched, and
correct every later bit --- a receiver tuned with mark and space swapped
decodes transparently, which real off-air reception needs.

\section{Verification, in layers}

The stack is exercised the way the rest of this book advocates. A native host
harness renders a transmission into a buffer and decodes it under injected
impairments: noise to the signal's own level, carrier offset to $\pm$80\,Hz,
baud error to $\pm$3\%, burst corruption, split-copy corruption, swapped
tones. The same sources then run on the board, where \texttt{decode air} and
\texttt{navtex air} drive the full-duplex codec path: the burst goes out the
speaker, crosses the room as sound, returns through the microphone, and
decodes character-perfect --- a genuinely analog channel with no cable.

Above the tests sits proof, using the pattern of Chapter~\ref{ch:spark}: the
protocol layers are SPARK, proved at silver (428 run-time checks, none
unproved) against the \texttt{FSK} spec as a contract boundary --- the Float
signal processing stays outside the surface, verified by the harness instead.
The prover earned its keep: it found a postcondition that was wrong for null
string slices, an out-of-range intermediate in two bit counters, and a
side-effecting function the subset rightly rejects. The example's
\texttt{prove/prove.sh} reruns the whole proof in one command.
